\SourcePage{590}{593}
\section{Nonspherical Nuclei}

A system of particles moving in a spherically symmetric field cannot have a rotational energy spectrum; in quantum mechanics, rotation has no meaning at all for such a system. This applies also to the nuclear shell model considered in the preceding section, with its spherically symmetric self-consistent field.

In quantum mechanics, the division of a system's energy into internal and rotational parts has no rigorous mean\SourcePageJoin{591}{594}ing in general. It can only be approximate, and is possible when, for one physical reason or another, the system can to good approximation be treated as an assembly of particles moving in a specified field that lacks spherical symmetry. The rotational structure of the levels then arises when the possibility that this field rotates relative to a fixed coordinate system is taken into account. We encountered such a case, for example, in molecules whose electronic terms can be defined as the energy levels of a system of electrons moving in the specified field of fixed nuclei.

Experiment shows that most nuclei do not, in fact, have rotational structure. This means that a spherically symmetric self-consistent field is a good approximation for them; in other words, the nuclei are spherical (to within quantum fluctuations).

There is, however, a class of nuclei whose energy spectra are of rotational type (including nuclei in the approximate ranges of atomic weights $150<A<190$ and $A>220$). For them, a spherically symmetric self-consistent field is an entirely unsuitable approximation. In principle, their self-consistent field must be sought without any prior assumptions about its symmetries, so that the nuclear shape too is determined in a ``self-consistent'' manner. Experiment shows that the correct model for this class is a self-consistent field possessing an axis of symmetry and a symmetry plane perpendicular to it (that is, the symmetry of an ellipsoid of revolution). The concept of nonspherical nuclei was developed most fully in the work of \emph{A.~Bohr and Mottelson} (A.~Bohr, B.~R.~Mottelson, 1952--1953).

We emphasize that these are two qualitatively different classes of nuclei. This is evident, in particular, from the fact that nuclei are either spherical or nonspherical with a degree of nonsphericity that is by no means small.

The presence of unfilled shells in a nucleus favors the emergence of nonsphericity; nucleon pairing apparently also plays an important part in this phenomenon. Conversely, shell closure favors a spherical nucleus. The doubly magic nucleus ${}^{208}_{82}\mathrm{Pb}$ is characteristic in this respect: because its nucleon configuration exhibits particularly strong closure, this nucleus (and nuclei close to it) is spherical, producing a break in the sequence of nonspherical heavy nuclei.

\SourcePage{592}{595}
The energy levels of a nonspherical nucleus are represented as the sum of two parts: the levels of the ``nonrotating'' nucleus and the energy of its rotation as a whole. In even-even nuclei, the intervals within the rotational structure of the levels are small compared with the separations between levels of the ``nonrotating'' nucleus.

The classification of the levels of a nonspherical nucleus is in many respects analogous to that of the levels of a diatomic molecule (composed of identical atoms), since in both cases the field in which the particles (nucleons or electrons) move has the same symmetry. We can therefore use directly a number of results obtained in Chapter~XI\footnote{We emphasize that the analogy concerns specifically the classification of the levels of a diatomic molecule, not of a symmetric top. For a system of particles moving in an axially symmetric field, rotation about the field axis has no meaning, just as rotation about any axis has no meaning for a system in a centrally symmetric field.}.

Consider first the classification of the states of the ``nonrotating nucleus.'' In an axially symmetric field, only the projection of angular momentum on the symmetry axis is conserved. Each state of the nucleus is therefore characterized first of all by the value $\Omega$ of the projection of its total angular momentum\footnote{Here $\Omega\geqslant0$ by definition, as is $\Lambda$ in diatomic molecules. Negative $\Omega$ could occur there only from $\Omega=\Lambda+\Sigma$, since $\Sigma$ has either sign according to the relative directions of orbital angular momentum and spin.}, which may take both integer and half-integer values. According to how the wave function behaves when the signs of the coordinates of all nucleons are changed (relative to the center of the nucleus), the levels are divided into even ($g$) and odd ($u$) ones.

In addition, when $\Omega=0$, positive and negative states are distinguished further, according to the behavior of the wave function under reflection in a plane through the nuclear axis (see \S~78).

The ground states of even-even nonspherical nuclei are $0_g$ states (the numeral indicates the value of $\Omega$), corresponding to zero angular momentum and to the highest symmetry of the wave function; this results from pairwise pairing of all neutrons and all protons. If instead the nucleus contains an odd number of protons or neutrons, the state of the ``odd'' nucleon in the self-consistent field of the even-even nuclear ``core'' may be considered.

The value of $\Omega$ is then determined by the projection of the angular momentum of this nucleon. Similarly, in an odd-odd nucleus, $\Omega$ is formed from the projections of the angular momenta of the odd neutron and the o\SourcePageJoin{593}{596}dd proton ($\Omega=|\omega_p\pm\omega_n|$).

At the same time, it must be emphasized that definite values cannot be assigned separately to the projections of the orbital angular momentum and the spin of the nucleon.

Although the spin-orbit coupling of the nucleon is small in comparison with the energy of its interaction with the self-consistent field of the core, it is not small compared with the separations between neighboring energy levels of the nucleon in this field. Yet it is precisely the latter condition that would be required for perturbation theory to apply and thus permit the orbital angular momentum and nucleon spin to be considered separately to a good approximation\footnote{In spherical nuclei, it was nevertheless possible to determine the value of $l$ by using conservation of parity and of angular momentum together.}.

We now turn to the rotational structure of the levels of a nonspherical nucleus. The intervals of this structure are small compared with the spin-orbit interaction of the nucleons in the nucleus; this situation corresponds to case $a$ in the theory of diatomic molecules (see \S~83).

The total angular momentum $\mathbf J$ of the rotating nucleus is of course conserved. For a given $\Omega$, its magnitude $J$ takes values beginning with $\Omega$:
\begin{equation}
 J=\Omega,\quad \Omega+1,\quad \Omega+2,\ldots
 \tag{119.1}
\end{equation}
(see (83.2)). There is a further restriction on the possible values of $J$ for nuclei with $\Omega=0$: in the states $0_g^+$ and $0_u^+$, $J$ takes only even values, whereas in the states $0_g^-$ and $0_u^-$ it takes odd values (see \S~86). In particular, in the rotational levels of the ground term of an even-even nucleus ($0_g^+$), $J$ takes the values $0,2,4,\ldots$

The rotational energy of the nucleus is given by
\begin{equation}
 E_{\mathrm{rot}}=\frac{\hbar^2}{2I}J(J+1),
 \tag{119.2}
\end{equation}
where $I$ is the moment of inertia of the nucleus (about an axis perpendicular to its symmetry axis); this formula corresponds to the analogous expression in the theory of diatomic molecules (the term depending

\SourcePage{594}{597}
on $J$ in (83.6)). The lowest level corresponds to the smallest possible value of $J$, namely $J=\Omega$.

By (119.2), the rotational structure of the levels is characterized by definite interval rules that, for a specified $\Omega$, do not depend on the other characteristics of the level. Thus, the components of the rotational structure of the ground term of an even-even nucleus (with $J=2,4,6,8,\ldots$) lie above the deepest level ($J=0$) by intervals in the ratios $1:3$, $3:7:12\ldots$

Formula (119.2), however, is insufficient for states with $\Omega=1/2$, which may occur in nuclei containing an odd number of nucleons.

In this case there is a contribution to the energy, comparable with (119.2), that is associated with the interaction of the odd nucleon with the centrifugal field of the rotating nucleus. Its dependence on $J$ can be found as follows.

As is known from mechanics (see I, \S~39), the energy of a particle in a rotating coordinate system contains an additional term equal to the product of the angular velocity of rotation and the angular momentum of the particle. The corresponding term in the nuclear Hamiltonian can be written as $2b\hat{\mathbf K}\hat{\boldsymbol\sigma}$, where $b$ is a constant, $\hat{\mathbf K}$ is the rotational angular momentum of the nuclear core (the nucleus without the last nucleon), and $\hat{\boldsymbol\sigma}$ is the angular momentum of the nucleon. The last quantity is to be understood here in a purely formal sense (in reality, an angular-momentum vector for a nucleon in the axial field of the nucleus does not exist): it is an operator analogous to the spin-$1/2$ operator, giving transitions between states whose angular-momentum projections are $\pm1/2$, in accordance with $\Omega=1/2$\footnote{The special feature of the case $\Omega=1/2$ is precisely the existence of matrix elements of the energy perturbation for transitions between states differing only in the sign of the angular-momentum projection and therefore belonging to the same energy. Consequently, an energy shift already appears in the first order of perturbation theory. This phenomenon is analogous to the $\Lambda$-doubling of the levels of a diatomic molecule with $\Omega=1/2$ (see \S~88).}. Since $\mathbf K=\mathbf J-\boldsymbol\sigma$, the eigenvalues of this operator are
\[
 2b\mathbf K\boldsymbol\sigma
 =b\left[J(J+1)-K(K+1)-\frac34\right].
\]
Adding, for convenience, the $J$-independent constant $b/2$, we find that this quantity is $\pm b(J+1/2)$ for $J=K\pm1/2$.

This expression can be written as $(-1)^{J-1/2}b(J+1/2)$ if we note that the angular momentum $K$ of the core (which is an even-even nucleus) is an even integer. Thus,

\SourcePage{595}{598}
we finally obtain the following expression for the rotational energy of a nucleus with $\Omega=1/2$:
\begin{equation}
 E_{\mathrm{rot}}=\frac{\hbar^2}{2I}J(J+1)
 +(-1)^{J-1/2}b\left(J+\frac12\right)
 \tag{119.3}
\end{equation}
(A.~Bohr, B.~Mottelson, 1953). Note that if the constant $b$ is positive and sufficiently large, the level with $J=3/2$ may lie below that with $J=1/2$; that is, the normal ordering of rotational levels, in which the lowest level corresponds to the smallest possible $J$, may be violated.

The moment of inertia of a nonspherical nucleus cannot be calculated as that of a rigid body having a specified shape. Such a calculation would be possible only if the nucleons moving in the self-consistent field of the nucleus could be regarded as not interacting directly with one another. In reality, the pairing phenomenon reduces the moment of inertia below the value corresponding to a rigid body.

The magnetic moment $\boldsymbol\mu$ of a nonspherical nucleus is the sum of the magnetic moment of the ``nonrotating'' nucleus and the moment associated with nuclear rotation. The first, after averaging over the motion of the nucleons in the nucleus, is directed along the nuclear axis. Denoting its magnitude by $\mu'$ and the unit vector along the nuclear axis by $\mathbf n$, we write it as $\mu'\mathbf n$. The magnetic moment associated with rotation, after the same averaging, is directed along $\mathbf J-\Omega\mathbf n$: the total mechanical angular momentum of the nucleus less the angular momentum of the nucleons in the ``nonrotating nucleus''\footnote{This form may be used only when $\Omega\ne1/2$ (see Problem~2).}.

Thus,
\begin{equation}
 \boldsymbol\mu=\mu'\mathbf n+g_r(\mathbf J-\Omega\mathbf n).
 \tag{119.4}
\end{equation}
Here $g_r$ is the gyromagnetic factor for nuclear rotation. Since only protons contribute to the magnetic moment during rotation,
\begin{equation}
 g_r=\frac{I_p}{I_p+I_n},
 \tag{119.5}
\end{equation}
where $I_n$ and $I_p$ are the neutron and proton parts of the nuclear moment of inertia (for a system consisting only of protons, one would simply have $g_r=1$). In general, the ratio (119.5) does not coincide with the ratio $Z/A$ of the number of protons to the total mass of the nucleus.

\SourcePage{596}{599}
After averaging over nuclear rotation, the magnetic moment is directed along the conserved vector $\mathbf J$:
\[
 \hat{\overline{\boldsymbol\mu}}
 =\frac\mu J\hat{\mathbf J}
 =(\mu'-\Omega g_r)\overline{\mathbf n}+g_r\hat{\mathbf J}.
\]
As usual, we multiply both sides of this equality by $\hat{\mathbf J}$ and pass to eigenvalues. For the ground state of the nucleus, $\Omega=J$, and we obtain
\begin{equation}
 \mu=(\mu'+g_r)\frac{J}{J+1}.
 \tag{119.6}
\end{equation}

\subsection*{Problems}

\ProblemHeading{1.} Express the quadrupole moment $Q$ of a rotating nucleus in terms of the quadrupole moment $Q_0$ referred to axes attached to the nucleus (A.~Bohr, 1951).

\SolutionHeading The operator of the quadrupole-moment tensor of a rotating nucleus is expressed in terms of $Q_0$ by
\[
 Q_{ik}=\frac32Q_0\left(n_in_k-\frac13\delta_{ik}\right);
\]
this is a symmetric traceless tensor formed from the components of the unit vector $\mathbf n$ along the nuclear axis, with $Q_{zz}=Q_0$. Averaging over the rotational state of the nucleus is carried out as in the solution of the problem in \S~29, except that $n_iJ_i=\Omega$ rather than zero, and gives an expression of the form (75.2), with
\[
 Q=Q_0\frac{3\Omega^2-J(J+1)}{(2J+3)(J+1)}.
\]
For the nuclear ground state, with $\Omega=J$, we obtain
\[
 Q=Q_0\frac{(2J-1)J}{(2J+3)(J+1)}.
\]
As $J$ increases, the ratio $Q/Q_0$ tends to 1, though rather slowly.

\ProblemHeading{2.} Determine the magnetic moment in the ground state of a nucleus with $\Omega=1/2$.

\SolutionHeading In this case the magnetic-moment operator can be written in terms of the operator $\hat{\boldsymbol\sigma}$ introduced in the text as
\[
 \hat{\boldsymbol\mu}=2\mu'\hat{\boldsymbol\sigma}+g_r\hat{\mathbf K},
 \qquad \hat{\mathbf K}=\hat{\mathbf J}-\hat{\boldsymbol\sigma}.
\]
The remainder of the calculation is analogous to that carried out in the text. If the ground level of the nucleus corresponds to $J=1/2$ (so that $K=J-1/2=0$), then $\mu=\mu'$. If instead $J=3/2$ in the ground state (so that $K=J+1/2=2$), then $\mu=\frac95g_r-\frac35\mu'$.

\ProblemHeading{3.} Determine the energies of the first few levels in the rotational structure of the ground state of an even-even nucleus having the symmetry of a triaxial ellipsoid.

\SourcePage{597}{600}
\SolutionHeading The ground state of an even-even nucleus corresponds to the most symmetric wave function of the ``nonrotating'' nucleus, namely a function whose symmetry belongs to the representation $A_1$ of the group $D_2$. For a specified value of $J$, there are therefore only $(J/2)+1$ distinct levels when $J$ is even, or $(J-1)/2$ when $J$ is odd. For $J=2$ they are given by formula (7) obtained in the problems to \S~103, and for $J=3$ by formula (8).
